\subsection{An evaluated elliptic Frobenius retaining the ten abelian coordinates}

\begin{statement}{E.58}{Proved}{Two actual elliptic quotients and their based Frobenius at two}
Use the actual degree-nine elliptic covering of E.50, with
$F=\pi_1^{(3)}(E\setminus\{O\},\vec w)$ and
$N=\ker(F\to F^{\rm ab}/3)$.
The specified cubic tangent $s=z/x$ maps to $\vec w$ under this
cover, by the exact morphism and tangent calculation below.
Choose the modular path with the principal compatible complex roots
of the Weierstrass discriminant, as specified in the proof.
Let $x,y$ be its free pro-three generators, put
\[
 c=[x,y],\quad d=[x,c],\quad e=[y,c],\qquad
 \Lambda=\mathbb Z_3[[X-1,Y-1]],\quad
 \mathfrak m=(X-1,Y-1),
 \tag{E58.1}
\]
and identify $F'/F''=\Lambda c$, with $X,Y$ acting by
conjugation by $x,y$. Write
\[
 U=X^3-1,\quad V=Y^3-1,\quad J=(U,V),\quad
 P=(1+X+X^2)(1+Y+Y^2),\quad I_B=J^3+PJ^2.
 \tag{E58.2}
\]
For an ideal $L\subset\Lambda$, let $\widetilde L$ denote the
inverse image of $Lc$ in $F'$. The following are actual quotient maps:
\[
 \begin{gathered}
 B=N/(F''\gamma_4N)=N/\widetilde I_B,
 \qquad R=N/\widetilde{J\mathfrak m},
 \qquad P_3=N/\gamma_4F,\\
 B\twoheadrightarrow R\twoheadrightarrow P_3,
 \qquad B\twoheadrightarrow R\twoheadrightarrow
 C_{\rm cc}=N/[N,F'].
 \end{gathered}
 \tag{E58.3}
\]
Here $C_{\rm cc}$ is exactly the cuspidally central quotient of E.36.
The comparison is the fiber product
\[
 R\ \simeq\ C_{\rm cc}\mathop{\times}_{Q_*}P_3,
 \qquad Q_*=N/\widetilde{J+\mathfrak m^2}.
 \tag{E58.4}
\]
The commutator-module quotients in $B,R,P_3,Q_*$ are, respectively,
\[
 \Lambda/I_B,\qquad
 M=\Lambda/(J\mathfrak m),\qquad
 \Lambda/\mathfrak m^2,\qquad
 \Lambda/(J+\mathfrak m^2)
       \simeq\mathbb Z_3\oplus(\mathbb Z/3)^2.
 \tag{E58.5}
\]
The first three are free of ranks $51,11,3$ over $\mathbb Z_3$.
Thus the kernels of $B\to R$, $B\to P_3$, $R\to P_3$ and
$R\to C_{\rm cc}$ are abelian and free of ranks $40,48,8,2$.
The abelianization of $R$ is the full $N^{\rm ab}$, of rank ten.

Use the following integral basis of $M$:
\[
 e_{ij}=X^iY^j\quad(0\le i,j\le2),\qquad U,V.
 \tag{E58.6}
\]
Its multiplication is determined by
\[
 \begin{split}
 e_{ij}e_{kl}
   &=e_{\overline{i+k},\overline{j+l}}
         +\left\lfloor\frac{i+k}{3}\right\rfloor U
         +\left\lfloor\frac{j+l}{3}\right\rfloor V,\\
 e_{ij}U&=U,\qquad e_{ij}V=V,
 \qquad U^2=UV=V^2=0,
 \end{split}
 \tag{E58.7}
\]
where the barred indices belong to $\{0,1,2\}$.
Every element of $R$ has unique coordinates
$\xi^a\eta^b c^H$, with
$\xi=x^3$, $\eta=y^3$, $a,b\in\mathbb Z_3$ and $H\in M$.
It has the exact lower-central lattices
\[
 \gamma_2R=\mathbb Z_3Pc\oplus\mathbb Z_3Uc\oplus\mathbb Z_3Vc,
 \qquad
 \gamma_3R=9\mathbb Z_3Uc\oplus9\mathbb Z_3Vc,
 \qquad\gamma_4R=1.
 \tag{E58.8}
\]
In particular $\gamma_2R/\gamma_3R\simeq
\mathbb Z_3\oplus(\mathbb Z/9)^2$. These finite summands remain
part of the integral quotient.

Choose the actual prolongation at the prime of
$K=\mathbb Q(\zeta_3)$ over two for which the compatible positive
roots of three lie in $\mathbb Q_2$. Let $\varphi$ be arithmetic
Frobenius there. Its genuine based action on $R$ is the following
fully evaluated automorphism. Put
\[
 \begin{split}
 H_x&=(1+X+X^2)(3+Y-Y^2)+\frac{31}{4}U+\frac32V,\\
 H_y&=(X^2-X-3)(1+Y+Y^2)-\frac32U-\frac{23}{4}V,\\
 \Phi(e_{ij})&=e_{ij}-iU-jV,
       \qquad\Phi(U)=-2U,\quad\Phi(V)=-2V.
 \end{split}
 \tag{E58.9}
\]
Then
\[
 \varphi(\xi)=\xi^{-2}c^{H_x},\qquad
 \varphi(\eta)=\eta^{-2}c^{H_y},\qquad
 \varphi(c^H)=c^{4\Phi(H)}.
 \tag{E58.10}
\]
Equivalently its map on all thirteen coordinates is
\[
 \begin{split}
 (a,b,H)\longmapsto\bigl(&-2a,-2b,
    4\Phi(H)+aH_x+bH_y+9a(a-1)U\\
    &\hspace{35mm}+[18ab-9b(b-1)]V\bigr).
 \end{split}
 \tag{E58.11}
\]
The induced ambient class-three action on $F/\gamma_4F$ is
\[
 \begin{array}{c|l}
 x&x^{-2}c^3d^{19/4}e^{1/2}\\
 y&y^{-2}c^{-3}d^{-1/2}e^{-11/4}\\
 c&c^4\\
 d&d^{-8}\\
 e&e^{-8}.
 \end{array}
 \tag{E58.12}
\]
Restriction to $P_3$ gives
\[
 \xi\longmapsto\xi^{-2}c^9d^{129/4}e^{3/2},\qquad
 \eta\longmapsto\eta^{-2}c^{-9}d^{-3/2}e^{-105/4},
 \tag{E58.13}
\]
with the last three images from (E58.12).

All these coordinates include an exact change from the
$\mathbb Q_2$ discriminant-root convention. Define
$T_1(a):x\mapsto xy^{-a},\ y\mapsto y$.
On $R$ and on $F/\gamma_4F$ the two actual coordinate actions satisfy
\[
 \varphi_{\rm principal}
    =T_1(1/2)\varphi_{\mathbb Q_2}T_1(-1/2)
    =\operatorname{Int}(c^{1/8})\varphi_{\mathbb Q_2}.
 \tag{E58.14}
\]
The second equality compares the two actions; the first is the
actual integral change of generators. In (E58.12), the
$\mathbb Q_2$ version replaces $19/4,-11/4$ by $9/2,-3$.
In (E58.9), it replaces $31/4,-23/4$ by $15/2,-6$.

The factorization proved below is through the actual image of
the pair consisting of the Tate representation and the specified
discriminant Kummer cocycle. It does not construct a section from
an unrestricted matrix group. Equations (E58.3)--(E58.5) retain the
remaining rank-forty kernel in the larger quotient $B$; its full
based Frobenius action is not evaluated here. The manuscript anchor
is \S51.16, scans 172--178, folios 733v--736v, equations
(371)--(379), with Part E(c),(d).
\end{statement}

\begin{proof}
\emph{The based elliptic morphism.}
On the compact cubic $x_0^3+y_0^3+z_0^3=0$, the formulas
\[
 X_0=-\frac{3z_0}{x_0+y_0},\qquad
 Y_0=\frac{9(x_0-y_0)}{x_0+y_0}
 \tag{E58.15}
\]
give $Y_0^2=4X_0^3-27$, with inverse
$(x_0:y_0:z_0)=(9+Y_0:9-Y_0:-6X_0)$.
They extend across the displayed poles as inverse morphisms of
the smooth compact curves, carrying $O=(1:-1:0)$ to the origin.
For $s=z_0/x_0$, the branch through $O$ has
$y_0/x_0=-(1+s^3)^{1/3}$. Consequently the Weierstrass parameter
is
\[
 t=-2X_0/Y_0
   =\frac{2s}{3(1+(1+s^3)^{1/3})}
   =s/3+O(s^4).
 \tag{E58.16}
\]
Multiplication by three has derivative three at the origin.
Thus the actual covering $[3]$ takes the original $s$-tangent to
the Weierstrass $t$-tangent of coefficient one. Its remaining
formal unit has constant term one, whose compatible roots also
have constant term one. This proves the based comparison used
here without introducing an additional tangent Kummer character.
The discriminant for $dX_0/Y_0$ is $\Delta=-3^9$.

\emph{The integral quotient rings and groups.}
The free metabelian module in E.50 gives $F'/F''=\Lambda c$.
In $N/F''$, the two handle generators are $x^3,y^3$, and the
whole module $\Lambda c$ is present. Their commutators give
$Pc,Uc,Vc$. Since $N$ is normal in $F$, its derived subgroup
is stable under $X,Y$; hence
\[
 N'/F''=(P,J)c,
 \qquad \gamma_rN\,F''/F''=J^{r-2}(P,J)c\quad(r\ge2).
 \tag{E58.17}
\]
This proves the ideal $I_B$ and every quotient in (E58.3).
Similarly $[N,F']/F''=Jc$ and $\gamma_4F/F''=\mathfrak m^2c$.

Here is the exact intersection, over $\mathbb Z_3$:
\[
                         J\cap\mathfrak m^2=J\mathfrak m.
 \tag{E58.18}
\]
Write $u=X-1$, $v=Y-1$.
For $Ua+Vb\in\mathfrak m^2$, its linear term is
$3a(0)u+3b(0)v$. Since $\mathbb Z_3$ has no three-torsion,
$a(0)=b(0)=0$, proving the nontrivial containment.
The usual explicit lifting argument for two normal quotients
now gives (E58.4), including surjectivity onto the fiber product.
The sum ideal gives $3u=3v=u^2=uv=v^2=0$, which is the last
module in (E58.5).

The ring defined by (E58.7) is free on its eleven displayed
vectors. Its multiplication is associative: the two integer carries
in each coordinate add independently. It satisfies $X^3=1+U$,
$Y^3=1+V$ and $\mathfrak m U=\mathfrak m V=0$.
Conversely these relations reduce every formal series to the
displayed basis. They define continuous maps in both directions
between this ring and $\Lambda/(J\mathfrak m)$, inverse on $X,Y$;
the powers of $\mathfrak m$ tend to zero three-adically in the
finite ring. This proves its integral normal form and topology.

For the rank in $B$, use $U,V$ as formal variables.
$\Lambda$ is free of rank nine over $\mathbb Z_3[[U,V]]$,
with basis $e_{ij}$, by distinguished-polynomial division.
Modulo $J^3$ it has rank $9(1+2+3)=54$.
The additional relations $PU^2,PUV,PV^2$ are independent
primitive vectors, one sum of all nine coset vectors in each
degree-two component. Multiplication by $X$ or $Y$ adds only
terms in $J^3$, so these are all the relations from $PJ^2$.
The rank is therefore $51$. The maps to the free modules
$M$ and $\Lambda/\mathfrak m^2$ are surjective, proving the
kernel ranks and their freeness. The same reasoning proves
the two smaller kernel ranks.

Modulo $J\mathfrak m$, the ideal $(P,J)$ is exactly
$\mathbb Z_3P+\mathbb Z_3U+\mathbb Z_3V$:
$(X-1)P=3U$, $(Y-1)P=3V$ there.
Furthermore $UP=9U$, $VP=9V$ and $J^2=0$.
These equalities prove (E58.8).
The three vectors $P,U,V$ are primitive and independent in
$M$, so the abelianization has rank $2+(11-3)=10$ with no
torsion. The subgroup $\widetilde{J\mathfrak m}$ lies in $N'$,
so its quotient identifies this abelianization with the actual
$N^{\rm ab}$.

\emph{The universal arithmetic factorization with its sign.}
We use the original modular-path construction of
\cite[\S\S3.6,5.9 and Theorem 6.10.3]{A_Nakamura2013}.
In its convention $z=[x,y]^{-1}=c^{-1}$.
On the kernel of the Tate representation, its equation (3.6.3)
identifies the action on $F/F''$ with conjugation by
$z^{-E_\sigma}=c^{E_\sigma}$.
The augmentation of this integral measure is
$\kappa_\Delta(\sigma)/12$.
In particular the kernel of the pair $(\rho_E,\kappa_\Delta)$
has $E_\sigma\in\mathfrak m$.
Conjugation by such a module element changes $x^3,y^3$ by
elements of $J\mathfrak m c$ and fixes $F'/F''$.
It consequently acts trivially on $R$.
It also acts trivially on $F/\gamma_4F$, since that commutator
change lies in $\mathfrak m^2c$.
This proves the stated factorization through the actual image
of the cocycle pair, with product
$(A,k)(B,l)=(AB,k+\det(A)l)$.
It specifies the sign of the cusp correction as well.

\emph{An explicit reference element.}
Nakamura's Weierstrass-section formulas
\cite[Proposition 5.7.3]{A_Nakamura2013} give the actual twists
\[
 T_1(a):x\mapsto xy^{-a},\ y\mapsto y,
 \qquad T_2(a):x\mapsto x,\ y\mapsto yx^a.
 \tag{E58.19}
\]
Their discriminant winding is $-a$. Direct multiplication gives
$(T_1(1)T_2(1))^6=\operatorname{Int}(c^{-1})$.
Define
\[
 D=T_2(-2)T_1(-1/2)T_2(-3)T_1(-1)T_2(-1).
 \tag{E58.20}
\]
It has Tate matrix $\operatorname{diag}(-2,-1/2)$ and
discriminant winding $15/2$.
These three-adic twist powers are actual continuous
automorphisms, and are images of the corresponding profinite
braid powers.

On the Tate tangential splitting, an actual Galois element whose
three-adic cyclotomic character is four has Tate matrix $\operatorname{diag}(1,4)$
and discriminant cocycle zero.
The explicit formula of \cite[Theorem 5.2.4]{A_Nakamura2013},
with the genuine coefficient $(4^2-1)/24=5/8$ verified in E.54,
has ambient class-three images
\[
 x\longmapsto x c^{3/2}d^{-3/2}e^{-5/4},
 \qquad y\longmapsto y^4.
 \tag{E58.21}
\]
There is also a convenient exact metabelian reference map
\[
 A_4:x\longmapsto c^{h(Y)}x,\ y\longmapsto y^4,
 \qquad h(Y)=\frac{Y^2+2Y+3}{(Y+1)(Y^2+1)}.
 \tag{E58.22}
\]
The denominator is a unit in $\Lambda$. Writing
$S_4=1+Y+Y^2+Y^3$, one has
$h=(S_4-4)/(Y^4-1)$, which proves directly that $A_4(c)=c^4$.
Its invertible Frattini map proves that it is an automorphism
of the free metabelian pro-three group.
Expansion to first order in $Y-1$ gives
$h=3/2-(5/4)(Y-1)+\cdots$, so it has exactly (E58.21).
The ratio of this map and the genuine Tate map is a
Tate-trivial, cusp-fixing automorphism; the regular-sequence
calculation in E.50 writes it as conjugation by $c^H$.
Equality of the class-three actions gives $H\in\mathfrak m$.
Thus (E58.22) and the genuine Tate map agree on $R$.
This last comparison is asserted on $R$ and on
$F/\gamma_4F$, with their displayed kernels.

The composite $DA_4$ therefore represents the actual pair
$(-2I,15/2)$ on these two quotients. Left composition by
$\operatorname{Int}(c^{-1/2})$ gives the pair
$(-2I,3/2)$. Left composition instead by
$\operatorname{Int}(c^{-5/8})$ gives $(-2I,0)$.
Collecting these explicit words proves (E58.9)--(E58.13).
For example the substitutions on the module are
$X\mapsto X^{-2}$, $Y\mapsto Y^{-2}$;
in (E58.7) these are precisely $\Phi$.
Powering the two handle images gives (E58.11).

\emph{The actual local Frobenius and the root transport.}
The good integral model is $Y_1^2+Y_1=X_0^3-7$.
It has three points over $\mathbb F_2$, hence Frobenius
polynomial $T^2+2$. The prime of $K$ over two has residue
field $\mathbb F_4$, so its Tate action is exactly $-2I$.
All three-power roots of three can be embedded compatibly
in $\mathbb Q_2$, where the odd-power map on units is bijective.
For every $n$, put
\[
 \Delta_n^{\rm loc}=-3^{9/3^n},\qquad
 \Delta_n^{\rm pr}=e^{\pi i/3^n}3^{9/3^n}.
 \quad
 \frac{\Delta_n^{\rm pr}}{\Delta_n^{\rm loc}}
       =\zeta_{3^n}^{(3^n+1)/2}.
 \tag{E58.23}
\]
These are compatible systems; the ratio has coherent exponent
$1/2\in\mathbb Z_3$. Frobenius fixes the first system and has
three-adic cyclotomic character four. Its discriminant cocycles are
therefore $0$ and $3/2$, respectively.

A modular path can be chosen to give the first system:
starting with any path, its root ratios give a compatible
three-adic exponent, and a twist (E58.19), whose discriminant
winding is surjective, corrects that exponent.
Changing that path by $T_1(1/2)$ changes the discriminant
cocycle by $(\chi-1)/2$, since its winding is $-1/2$.
Its matrix commutes with the scalar $-2I$, so the resulting
local Tate matrix remains exactly the one used above.
The actual change of coordinates is the first equality in
(E58.14); direct collection verifies it on the two ambient
class-three generators and on every generator of $R$.
The second equality follows by the computed difference
$\kappa_\Delta/12=(3/2)/12=1/8$.
Every exponent used in this transport is integral over
$\mathbb Z_3$.

The retained programs verify the cubic morphism, its tangent,
the point count, the exact braid, Tate, and root-change formulas,
and all thirteen generator images. They check the finite ring's
1,331 associativity identities, 144 exact group/action
comparisons and 156 power identities. The projection to the
ambient three-coordinate commutator module has an explicit
integer right inverse. These checks accompany the integral
ideal and group proofs above.
\end{proof}
